\documentclass[11pt]{article}
% Conference-style write-up skeleton for an RL / robot learning experiment:
% abstract, introduction, related work, setup, method, experiments, conclusions;
% theory organized as definitions, assumptions, theorems; link the code.
\usepackage[margin=1in]{geometry}
\usepackage{amsmath,amssymb,amsthm}
\usepackage{graphicx,booktabs,natbib,hyperref}
\usepackage[ruled,vlined]{algorithm2e}

\newtheorem{theorem}{Theorem}
\newtheorem{lemma}[theorem]{Lemma}
\newtheorem{proposition}[theorem]{Proposition}
\theoremstyle{definition}
\newtheorem{definition}{Definition}
\newtheorem{assumption}{Assumption}
\theoremstyle{remark}
\newtheorem{remark}{Remark}

\newcommand{\E}{\mathbb{E}}
\newcommand{\cS}{\mathcal{S}}
\newcommand{\cA}{\mathcal{A}}
\newcommand{\cO}{\mathcal{O}}

\title{TITLE}
\author{AUTHORS}
\date{}

\begin{document}
\maketitle

\begin{abstract}
Problem, why it matters, what you did, the main result in numbers, and what it implies. Four to six sentences.
\end{abstract}

\section{Introduction}
Motivation, the question, contributions as a short list, and the paper's organization.

\section{Related work}
Group by idea (access models, structural conditions, practical reset-based exploration), not by paper.

\section{Problem setup}
\begin{definition}[Access protocol]
State exactly what the learner may query.
\end{definition}
\begin{assumption}[Structure]
State each assumption once; refer to it by name later.
\end{assumption}

\section{Method}
\begin{algorithm}[t]
\caption{METHOD}
\KwIn{...}
\KwOut{...}
...
\end{algorithm}
List every choice that departs from the source paper (constants, function class, reset semantics).

\section{Experiments}
Environments, the swept variable, baselines, the fixed evaluator and its metrics, batch arithmetic and compute, number of seeds, how intervals are computed (IQM with stratified bootstrap, 95\%).
\begin{table}[t]\centering
\begin{tabular}{lcc}\toprule
Method & IQM return & 95\% CI \\ \midrule
 & & \\ \bottomrule
\end{tabular}
\caption{...}
\end{table}

\section{Discussion and conclusions}
What held, what did not and why, limitations, next steps.

\paragraph{Code.} \url{REPO-URL}

\bibliographystyle{plainnat}
\bibliography{refs}
\end{document}
